\documentclass[12pt]{article}
\usepackage{amsmath}
\usepackage{amssymb}
\usepackage{amsfonts}
\usepackage{graphicx}
\usepackage[hidelinks]{hyperref}
\usepackage{geometry}
\usepackage{multicol}

\geometry{
 a4paper,
 left=3cm,
 top=2cm,
 right=2cm,
 bottom=2cm
}

\begin{document}

\begin{center}
    \textbf{\large UNIVERSIDADE FEDERAL DO RIO GRANDE DO SUL} \\
    \textbf{\large INSTITUTO DE MATEMÁTICA E ESTATÍSTICA} \\
    \vspace{3cm}
    \textbf{\Large A FIRST COURSE ON DIFFERENTIAL EQUATIONS WITH MODELING APPLICATION BY DENNIS G. ZILL (10e)} \\
    \vspace{1cm}
    \textbf{\large LIST OF RESULTS} \\
    \vspace{3cm}
    
    \author{Leonardo Suarez}
    
    \vfill
    \textbf{\large Porto Alegre, RS} \\
    \textbf{\large\today} \\
\end{center}\newpage

\section*{List of Results}

\noindent Theorem 1.2.1 (\emph{Existence of a Unique Solution}) \dotfill 17 \\
Theorem 2.4.1 (\emph{Criterion for an Exact Differential}) \dotfill 66 \\
Theorem 4.1.1 (\emph{Existence of a Unique Solution}) \dotfill 119 \\
Theorem 4.1.2 (\emph{Superposition Principle—Homogeneous Equations}) \dotfill 122 \\
Theorem 4.1.3 (\emph{Criterion for Linearly Independent Solutions}) \dotfill 124 \\
Theorem 4.1.4 (\emph{Existence of a Fundamental Set}) \dotfill 125 \\
Theorem 4.1.5 (\emph{General Solution—Homogeneous Equations}) \dotfill 126 \\
Theorem 4.1.6 (\emph{General Solution—Nonhomogeneous Equations}) \dotfill 127 \\
Theorem 4.1.7 (\emph{Superposition Principle—Nonhomogeneous Equations}) \dotfill 128 \\
Theorem 4.2.1 (\emph{Reduction of Order}) \dotfill 133 \\

\vspace{0.5em}

\noindent Theorem 7.1.1 (\emph{Transforms of Some Basic Functions}) \dotfill 289 \\
Theorem 7.1.2 (\emph{Sufficient Conditions for Existence}) \dotfill 290 \\
Theorem 7.1.3 (\emph{Behavior of \(F(s)\) as \(s \to \infty\)}) \dotfill 291 \\
Theorem 7.2.1 (\emph{Some Inverse Transforms}) \dotfill 294 \\
Theorem 7.2.2 (\emph{Transform of a Derivative}) \dotfill 296 \\
Theorem 7.3.1 (\emph{First Translation Theorem}) \dotfill 302 \\
Theorem 7.3.2 (\emph{Second Translation Theorem}) \dotfill 306 \\
Theorem 7.4.1 (\emph{Derivatives of Transforms}) \dotfill 315 \\
Theorem 7.4.2 (\emph{Convolution Theorem}) \dotfill 318 \\
Theorem 7.4.3 (\emph{Transform of a Periodic Function}) \dotfill 323 \\
Theorem 7.5.1 (\emph{Transform of the Dirac Delta Function}) \dotfill 329 \\

\vspace{0.5em}

\noindent Theoreom 8.1.1 (\emph{Existence of a Unique Solution}) \dotfill 349 \\
Theorem 8.1.2 (\emph{Superposition Principle}) \dotfill 349 \\
Theorem 8.1.3 (\emph{Criterion for Linearly Independent Solutions}) \dotfill 350 \\
Theorem 8.1.4 (\emph{Existence of a Fundamental Set}) \dotfill 351 \\
Theorem 8.1.5 (\emph{General Solution—Homogeneous Systems}) \dotfill 351 \\
Theorem 8.1.6 (\emph{General Solution—Nonhomogeneous Systems}) \dotfill 352 \\
Theorem 8.2.1 (\emph{Distinct Real Eigenvalues: General Solution—Homogeneous Systems}) \dotfill 355 \\
Theorem 8.2.2 (\emph{Solutions Corresponding to a Complex Eigenvalue}) \dotfill 363 \\
Theorem 8.2.3 (\emph{Real Solutions Corresponding to a Complex Eigenvalue}) \dotfill 364 \\\newpage

\section*{1 Introduction to Differential Equations}

\noindent\textbf{Theorem 1.2.1} (\emph{Existence of a Unique Solution}) Let \(R\) be a rectangular region in the \(xy\)-plane defined by \(a \leq x \leq b\), \(c \leq y \leq d\) that contains the point \(\left(x_0, y_0\right)\) in its interior. If \(f(x,y)\) and \(\partial{f}/\partial{y}\) are continuous on \(R\), then there exists some interval \(I_0: (x_0 - h, x_0 + h)\), \(h > 0\), contained in \([a, b]\), and a unique function \(y(x)\) defined on \(I_0\), that is a solution of the initial-value problem. 

\section*{2 First-Order Differential Equations}

\noindent\textbf{Theorem 2.4.1} (\emph{Criterion for an Exact Differential}) Let \(M(x,y)\) and \(N(x,y)\) be continuous and have continuous partial derivatives on a rectangle \(R\) in the \(xy\)-plane. Then a necessary and sufficient condition that \(M(x,y)dx + N(x,y)dy\) be an exact differential is

\begin{equation}
    \frac{\partial M}{\partial y} = \frac{\partial N}{\partial x}
\end{equation}

\section*{4 Higher-Order Differential Equations}

\noindent\textbf{Theorem 4.1.1} (\emph{Existence of a Unique Solution}) Let \(a_n(x), a_{n-1}(x), \ldots, a_0(x)\) and \(g(x)\) be continuous on an interval \(I\) and let \(a_n(x)\neq 0\) for every \(x\) in \(I\). If \(x = x_0\) is any point in \(I\), then a solution \(y(x)\) of the initial-value problem exists on \(I\) and is unique.\\

\noindent\textbf{Theorem 4.1.2} (\emph{Superposition Principle—Homogeneous Equations}) Let \(y_1(x), y_2, \ldots, y_k(x)\) be solutions of a homogeneous \(n\)-th order linear differential equation. Then the linear combination

\begin{equation}
    y(x) = c_1y_1(x) + c_2y_2(x) + \cdots + c_ky_k(x),
\end{equation}

\noindent where \(c_i\), \(i = 1, 2, \ldots, k\) are arbitrary constants, is also a solution of the equation.\\

\noindent\textbf{Theorem 4.1.3} (\emph{Criterion for Linearly Independent Solutions}) Let \(y_1(x), y_2(x), \ldots, y_n(x)\) be \(n\) solutions of a homogeneous \(n\)-th order linear differential equation on an interval \(I\). Then the set of solutions is linearly independent on \(I\) iff the Wronskian \(W(y_1, y_2, \ldots, y_n) \neq 0\) for every \(x\) in \(I\), where

\begin{equation}
    W(y_1, y_2, \ldots, y_n) = \begin{vmatrix}
        y_1 & y_2 & \cdots & y_n\\
        y_1' & y_2' & \cdots & y_n'\\
        \vdots & & & \vdots\\
        y_1^{(n-1)} & y_2^{(n-1)} & \cdots & y_n^{(n-1)}
    \end{vmatrix}.
\end{equation}\\

\noindent\textbf{Theorem 4.1.4} (\emph{Existence of a Fundamental Set}) There exists a fundamental set of solutions for the homogeneous \(n\)-th order linear differential equation on an interval \(I\).\\

\noindent\textbf{Theorem 4.1.5} (\emph{General Solution—Homogeneous Equations}) Let \(y_1(x), y_2(x), \ldots, y_n(x)\) be a fundamental set of solutions of a homogeneous \(n\)-th order linear differential equation on an interval \(I\). Then the \emph{general solution} of the equation on the interval is given by

\begin{equation}
    y = c_1y_1(x) + c_2y_2(x) + \cdots + c_ny_n(x),
\end{equation}

\noindent where the \(c_i\), \(i = 1, 2, \ldots, n\) are arbitrary constants.\\

\noindent\textbf{Theorem 4.1.6} (\emph{General Solution—Nonhomogeneous Equations}) Let \(y_p(x)\) be any particular solution of a nonhomogeneous \(n\)-th order linear differential equation on an interval \(I\) and let \(y_1, y_2, \ldots, y_n\) be a fundamental set of solutions of the associated homogeneous equation on the same interval. Then the \emph{general solution} of the nonhomogeneous equation on the interval is given by

\begin{equation}
    y = c_1y_1(x) + c_2y_2(x) + \cdots + c_ny_n(x) + y_p(x),
\end{equation}

\noindent where the \(c_i\), \(i = 1, 2, \ldots, n\) are arbitrary constants.\\

\noindent\textbf{Theorem 4.1.7} (\emph{Superposition Principle—Nonhomogeneous Equations}) Let \(y_{p1}, y_{p2}, \ldots, y_{pk}\) be \(k\) particular solutions of a nonhomogeneous \(n\)-th order linear differential equation on an interval \(I\) corresponding, in turn, to \(k\) distinct functions \(g_1(x), g_2(x), \ldots, g_k(x)\). That is, suppose \(y_{pi}\) denotes a particular solution of the corresponding differential equation

\begin{equation}
    a_n(x)y^{(n)} + a_{n-1}(x)y^{(n-1)} + \cdots + a_0(x)y = g_i(x), \quad i = 1, 2, \ldots, k.
\end{equation}\\

\noindent then

\begin{equation}
    y_p(x) = y_{p1}(x) + y_{p2}(x) + \cdots + y_{pk}(x)
\end{equation}

\noindent is a particular solution of

\begin{equation}
    a_n(x)y^{(n)} + a_{n-1}(x)y^{(n-1)} + \cdots + a_0(x)y = g_1(x) + g_2(x) + \cdots + g_k(x).
\end{equation}

\noindent\textbf{Theorem 4.2.1} (\emph{Reduction of Order}) Let \(y_1(x)\) be a solution of the homogeneous \(n\)-th order linear differential equation on an interval \(I\) and that \(y_1(x) \neq 0\) for all \(x\) in \(I\). Then

\begin{equation}
    y_2(x) = y_1(x) \int \frac{e^{-\int P(x)dx}}{y_1^2(x)}dx
\end{equation}

\noindent is a second solution of the equation on the interval.\\

\newpage

\section*{7 The Laplace Transform}

\noindent\textbf{Definition 7.1.1} (\emph{Laplace Transform}) Let \(f\) be a function defined on \([0,\infty)\). Then the function \(F\) defined by

\begin{equation}
    \mathcal{L} \{f(t)\} = F(s) = \int_0^\infty e^{-st} f(t) dt
\end{equation}

\noindent is said to be the \emph{Laplace transform} of \(f\). The domain of \(F\) is the set of all \(s\) for which the improper integral converges. \\

\noindent\textbf{Theorem 7.1.1} (\emph{Transforms of Some Basic Functions})

\begin{multicols}{2}
    \begin{itemize}
        \item [(a)] \(\displaystyle\mathcal{L} \{1\} = \frac{1}{s}\)
        \item [(b)] \(\displaystyle\mathcal{L} \{t^n\} = \frac{n!}{s^{n+1}}\) for \(n = 0, 1, 2, \ldots\)
        \item [(c)] \(\displaystyle\mathcal{L} \{e^{at}\} = \frac{1}{s-a}\)
        \item [(d)] \(\displaystyle\mathcal{L} \{\sin(\omega t)\} = \frac{\omega}{s^2 + \omega^2}\)
        \item [(e)] \(\displaystyle\mathcal{L} \{\cos(\omega t)\} = \frac{s}{s^2 + \omega^2}\)
        \item [(f)] \(\displaystyle\mathcal{L} \{\sinh(\omega t)\} = \frac{\omega}{s^2 - \omega^2}\)
        \item [(g)] \(\displaystyle\mathcal{L} \{\cosh(\omega t)\} = \frac{s}{s^2 - \omega^2}\)
    \end{itemize}
\end{multicols}

\noindent\textbf{Definition 7.1.2} (\emph{Exponential Order}) A function \(f\) defined on \([0, \infty)\) is said to be of \emph{exponential order} as \(t\to\infty\) if there exist constants \(c\), \(M > 0\), and \(T > 0\) such that \(\vert f(t) \vert \leq M e^{ct}\) for all \(t > T\). If \(f\) is of exponential order, then the Laplace transform \(\mathcal{L} \{f(t)\} \leq Me^{ct}\) for all \(t > T\).\\

\noindent\textbf{Theorem 7.1.2} (\emph{Sufficient Conditions for Existence}) If \(f\) is piecewise continuous on \([0, \infty)\) and of exponential order as \(t \to \infty\), then the Laplace transform \(\mathcal{L} \{f(t)\}\) exists for all \(s > c\), where \(c\) is a constant such that \(\vert f(t) \vert \leq Me^{ct}\) for all \(t > T\).\\

\noindent\textbf{Theorem 7.1.3} (\emph{Behavior of \(F(s)\) as \(s \to \infty\)}) If \(f\) is piecewise continuous on \([0, \infty)\) and of exponential order as \(t \to \infty\) and \(\mathcal{L} \{f(t)\} = F(s)\), then \(\displaystyle\lim_{s \to \infty} F(s) = 0\).\\

\noindent\textbf{Theorem 7.2.1} (\emph{Some Inverse Transforms})

\begin{multicols}{2}
    \begin{itemize}
        \item [(a)] \(\displaystyle\mathcal{L}^{-1} \left\{ \frac{1}{s} \right\} = 1\)
        \item [(b)] \(\displaystyle\mathcal{L}^{-1} \left\{ \frac{n!}{s^{n+1}} \right\} = t^n\) for \(n = 0, 1, 2, \ldots\)
        \item [(c)] \(\displaystyle\mathcal{L}^{-1} \left\{ \frac{1}{s-a} \right\} = e^{at}\)
        \item [(d)] \(\displaystyle\mathcal{L}^{-1} \left\{ \frac{\omega}{s^2 + \omega^2} \right\} = \sin(\omega t)\)
        \item [(e)] \(\displaystyle\mathcal{L}^{-1} \left\{ \frac{s}{s^2 + \omega^2} \right\} = \cos(\omega t)\)
        \item [(f)] \(\displaystyle\mathcal{L}^{-1} \left\{ \frac{\omega}{s^2 - \omega^2} \right\} = \sinh(\omega t)\)
        \item [(g)] \(\displaystyle\mathcal{L}^{-1} \left\{ \frac{s}{s^2 - \omega^2} \right\} = cosh(\omega t)\)
    \end{itemize}
\end{multicols}

\noindent\textbf{Theorem 7.2.2} (\emph{Transform of a Derivative}) Suppose \(f\), \(f'\), \ldots, \(f^{n-1}\) are continuous on \([0, \infty)\), and \(f\), \(f'\), \ldots, \(f^{n-1}\) are of exponential order as \(t \to \infty\). Then for \(s > c\),

\begin{equation}
    \mathcal{L} \{f^{(n)}(t)\} = s^n F(s) - s^{n-1} f(0) - s^{n-2} f'(0) - \ldots - f^{(n-1)}(0)
\end{equation}

\noindent where \(F(s) = \mathcal{L} \{f(t)\}\).\\

\noindent\textbf{Theorem 7.3.1} (\emph{First Translation Theorem}) If \(\mathcal{L} \{f(t)\} = F(s)\) exists for \(s > c\) and \(a\) is any real rumber, then

\begin{equation}
    \mathcal{L} \{e^{at} f(t)\} = F(s-a)
\end{equation}

\noindent for \(s > c + a\).\\

\noindent\textbf{Theorem 7.3.2} (\emph{Second Translation Theorem}) If \(\mathcal{L} \{f(t)\} = F(s)\) exists for \(s > c \geq 0\), and \(a\) is a positive constant, then

\begin{equation}
    \mathcal{L} \{f(t-a)\ \mathcal{U}(t-a)\} = e^{-as}F(s)
\end{equation}

\noindent for \(s > c\).\\

\noindent\textbf{Theorem 7.4.1} (\emph{Derivatives of Transforms}) If the function \(f\) is piecewise continuous on \([0,\infty)\), of exponential order and \(\mathcal{L} \{f(t)\} = F(s)\), then for \(n = 1, 2, 3, \ldots\) and \(s > c\),

\begin{equation}
    \mathcal{L} \{t^n f(t)\} = (-1)^n\frac{d^n}{ds^n} F(s).
\end{equation}

\noindent\textbf{Theorem 7.4.2} (\emph{Convolution Theorem}) If the functions \(f\) and \(g\) are piecewise continuous on \([0,\infty)\), of exponential order and \(\mathcal{L} \{f(t)\} = F(s)\), \(\mathcal{L} \{g(t)\} = G(s)\), then for \(s > c\),

\begin{equation}
    \mathcal{L} \{f*g\} = \mathcal{L} \{f(t)\} \mathcal{L} \{g(t)\} = F(s)G(s).
\end{equation}

\noindent\textbf{Theorem 7.4.3} (\emph{Transform of a Periodic Function}) If \(f\) is piecewise continuous \([0,\infty)\), of exponential order, and periodic with period \(T\), then

\begin{equation}
    \mathcal{L} \{f(t)\} = \frac{1}{1-e^{-sT}} \int_0^T e^{-st} f(t) dt.
\end{equation}

\noindent\textbf{Theorem 7.5.1} (\emph{Transform of the Dirac Delta Function}) For \(t_0 > 0\),

\begin{equation}
    \mathcal{L} \{\delta(t - t_0)\} = e^{-st_0}.
\end{equation}

\newpage

\section*{8 Systems of Linear Differential Equations}

\noindent\textbf{Definition 8.1.1} (\emph{Solution Vector}) A \textbf{solution vector} on an interval \(I\) is any column matrix

\begin{equation}
    \mathbf{X} = \begin{pmatrix}
        x_1(t) \\
        x_2(t) \\
        \vdots \\
        x_n(t)
    \end{pmatrix}
\end{equation}

\noindent whose entries are differentiable functions satisfying the system \(\mathbf{X'} = \mathbf{AX} + \mathbf{F}\) on the interval.\\

\noindent\textbf{Theorem 8.1.1} (\emph{Existence of a Unique Solution}) Let the entries of the matrices \(\mathbf{A}(t)\) and \(\mathbf{F}(t)\) be functions continuous on a common interval \(I\) that contains the point \(t_0\). Then there exists a unique solution of the initial-value problem \(\mathbf{X'} = \mathbf{A}(t)\mathbf{X} + \mathbf{F}(t)\), \(\mathbf{X}(t_0) = \mathbf{X_0}\) on the interval.\\

\noindent\textbf{Theorem 8.1.2} (\emph{Superposition Principle}) Let \(\mathbf{X}_1\), \(\mathbf{X}_2\), \(\ldots\), \(\mathbf{X}_k\) be a set of solution vectors of the homogeneous system \(\mathbf{X'} = \mathbf{A}\mathbf{X}\) on an interval \(I\). Then the linear combination

\begin{equation}
    \mathbf{X} = c_1\mathbf{X}_1 + c_2\mathbf{X}_2 + \cdots + c_k\mathbf{X}_k,
\end{equation}

\noindent where \(c_i\), \(i = 1, 2, \ldots, k\) are arbitrary constants, is also a solution of the system on the interval.\\

\noindent\textbf{Theorem 8.1.3} (\emph{Criterion for Linearly Independent Solutions}) Let

\begin{equation}
    \mathbf{X}_1 = \begin{pmatrix}
        x_{11}\\
        x{21}\\
        \vdots\\
        x_{n1}
    \end{pmatrix}, \quad
    \mathbf{X}_2 = \begin{pmatrix}
        x_{12}\\
        x_{22}\\
        \vdots\\
        x_{n2}
    \end{pmatrix}, \quad \ldots, \quad
    \mathbf{X}_n = \begin{pmatrix}
        x_{1n}\\
        x_{2n}\\
        \vdots\\
        x_{nn}
    \end{pmatrix}
\end{equation}

\noindent be \(n\) solution vectors of an homogeneous system on an interval \(I\). Then the set of solution vectors is linearly independent on \(I\) iff the \textbf{Wronskian}

\begin{equation}
    W(\mathbf{X}_1, \mathbf{X}_2, \ldots, \mathbf{X}_n) = \begin{vmatrix}
        x_{11} & x_{12} & \cdots & x_{1n}\\
        x_{21} & x_{22} & \cdots & x_{2n}\\
        \vdots & & & \vdots\\
        x_{n1} & x_{n2} & \cdots & x_{nn} 
    \end{vmatrix} \neq 0
\end{equation}

\noindent for every \(t\) in the interval.\\

\noindent\textbf{Theorem 8.1.4} (\emph{Existence of a Fundamental Set}) There exists a fundamental set of solutions for the homogeneous system \(\mathbf{X'} = \mathbf{AX}\) on an interval \(I\).\\

\noindent\textbf{Theorem 8.1.5} (\emph{General Solution—Homogeneous Systems}) Let \(\mathbf{X}_1\), \(\mathbf{X}_2\), \(\ldots\), \(\mathbf{X}_n\) be a fundamental set of solutions of a homogeneous system on an interval \(I\). Then the \textbf{general solution} of the system on the interval is

\begin{equation}
    \mathbf{X} = c_1\mathbf{X}_1 + c_2\mathbf{X}_2 + \cdots + c_n\mathbf{X}_n,
\end{equation}

\noindent where the \(c_i\), \(i = 1, 2, \ldots, n\) are arbitrary constants.\\

\noindent\textbf{Theorem 8.1.6} (\emph{General Solution—Nonhomogeneous Systems}) Let \(\mathbf{X}_p\) be a given solution of a nonhomogeneous system on an interval \(I\) and let

\begin{equation}
    \mathbf{X}_c = c_1\mathbf{X}_1 + c_2\mathbf{X}_2 + \cdots + c_n\mathbf{X}_n
\end{equation}

\noindent denote the general solution on the same interval of the associated homogeneous system. Then the \textbf{general solution} of the nonhomogeneous system on the interval is

\begin{equation}
    \mathbf{X} = \mathbf{X}_c + \mathbf{X}_p.
\end{equation}

\noindent The general solution \(\mathbf{X}_c\) of the associated homogeneous system is called the \textbf{complementary function} of the nonhomogeneous system.\\

\noindent\textbf{Theorem 8.2.1} (\emph{Distinct Real Eigenvalues: General Solution—Homogeneous Systems}) Let \(\lambda_1\), \(\lambda_2\), \(\ldots\), \(\lambda_n\) be \(n\) distinct real eigenvalues of the coefficient matrix \(\mathbf{A}\) of the homogeneous system \(\mathbf{X'} = \mathbf{AX}\) and let \(\mathbf{K}_1\), \(\mathbf{K}_2\), \(\ldots\), \(\mathbf{K}_n\) be the corresponding eigenvectors. Then the \textbf{general solution} of the system on the interval \((-\infty, \infty)\) is given by

\begin{equation}
    \mathbf{X} = c_1\mathbf{K}_1 e^{\lambda_1}t + c_2\mathbf{K}_2 e^{\lambda_2}t + \cdots + c_n\mathbf{K}_n e^{\lambda_n}t.
\end{equation}

\noindent\textbf{Theorem 8.2.2} (\emph{Solutions Corresponding to a Complex Eigenvalue}) Let \(A\) be the coefficient matrix having real entries of the homogeneous system \(\mathbf{X'} = \mathbf{AX}\) and let \(\mathbf{K}_1\) be an eigenvector corresponding to the complex eigenvalue \(\lambda_1 = \alpha + i\beta\), \(\alpha\) and \(\beta\) real numbers. Then

\begin{equation}
    \mathbf{K}_1 e^{\lambda_1 t} \quad \text{and} \quad \overline{\mathbf{K}_1} e^{\overline{\lambda_1} t}
\end{equation}

\noindent are solutions of the system, where \(\overline{\mathbf{K}_1}\) is the complex conjugate of \(\mathbf{K}_1\) and \(\overline{\lambda_1} = \alpha - i\beta\).\\

\noindent\textbf{Theorem 8.2.3} (\emph{Real Solutions Corresponding to a Complex Eigenvalue}) Let \(\lambda_1 = \alpha + i\beta\) be a complex eigenvalue of the coefficient matrix \(\mathbf{A}\) of the homogeneous system \(\mathbf{X'} = \mathbf{AX}\) and let \(\mathbf{B}_1\) and \(\mathbf{B}_2\) denote the column vectors defined as \(\mathbf{B}_1 = \frac{1}{2} \left(\mathbf{K}_1 + \overline{\mathbf{K}_1}\right)\) and \(\mathbf{B}_2 = \frac{i}{2} \left(-\mathbf{K}_1 + \overline{\mathbf{K}_1}\right)\). Then

\begin{align}
    \mathbf{X}_1 &= \mathbf{B}_1 e^{\alpha t} \cos(\beta t) + \mathbf{B}_2 e^{\alpha t} \sin(\beta t)\\
    \mathbf{X}_2 &= -\mathbf{B}_2 e^{\alpha t} \cos(\beta t) + \mathbf{B}_1 e^{\alpha t} \sin(\beta t)
\end{align}

\noindent are linearly independent solutions of the system \(\mathbf{X'} = \mathbf{AX}\) on the interval \((-\infty, \infty)\).

\newpage

The following is an introduction to difference equations, based on \emph{An Introduction to Difference Equations} (3e), by S. Elaydi.

\section*{1 Dynamics of First-Order Difference Equations}

A linear \emph{homogeneous} first-order difference equation is given by

\begin{equation}
    x(n+1) = a(n)x(n), \quad x(n_0) = x_0, \quad n \geq n_0 \geq 0, \label{eq:difference_homo}
\end{equation}

\noindent and the associated \emph{nonhomogeneous} equation is

\begin{equation}
    y(n+1) = a(n)y(n) + g(n), \quad y(n_0) = y_0, \quad n \geq n_0 \geq 0, \label{eq:difference_nonhomo}
\end{equation}

\noindent where in both cases \(a(n) \neq 0\), and \(a(n)\) and \(g(n)\) are real-valued functions defined for \(n \geq n_0\).

The solution to \eqref{eq:difference_homo} can be obtained by iteration and mathematical induction and is given by

\begin{equation}
    x(n) = \left[\prod_{i=n_0}^{n-1}a(i)\right] x_0
\end{equation}

For the nonhomogeneous case, the unique solution is

\begin{equation}
    y(n) = \left[\prod_{i=n_0}^{n-1}a(i)\right]y_0 +  \sum_{r=n_0}^{n-1}\left[\prod_{i=r+1}^{n-1}a(i)\right]g(r).
\end{equation}

\noindent\textbf{Definition 1.5} A point \(x^*\) in the domain of \(f\) is said to be an \emph{equilibrium point} of a difference equation 

\begin{equation}
    x(n+1) = f(x(n))\label{eq:difference_eq}
\end{equation} 

\noindent if it is a fixed point of \(f\), i.e., \(f(x^*) = x^*\).\\

\noindent\textbf{Definition 1.6} Let \(x\) be a point in the domain of \(f\). If there exists a positive integer \(r\) and an equilibrium point \(x^*\) of \eqref{eq:difference_eq} such that \(f^r(x) = x^*\), \(f^{r-1} \neq x^*\), then \(x\) is an \emph{eventually equilibrium (fixed) point}.\\

\noindent\textbf{Definition 1.8} 

\begin{itemize}
    \item [(a)] The equilibrium point \(x^*\) of \eqref{eq:difference_eq} is \emph{stable} if given \(\epsilon > 0\) there exists \(\delta > 0\) such that \(\vert x_0 - x^* \vert < \delta\) implies \(\vert f^n(x_0) - x^* \vert < \epsilon\) for all \(n > 0\). If \(x^*\) is not stable, then it is called \emph{unstable}.
    \item [(b)] The point \(x^*\) is said to be \emph{attracting} if there exists \(\eta > 0\) such that \[\vert x(0) - x^* \vert < \eta \implies \lim_{n\to\infty} x(n) = x^*.\] If \(\eta = \infty\), \(x^*\) is called a \emph{global attractor} or \emph{globally attracting}.
    \item [(c)] The point \(x^*\) is an \emph{asymptotically stable equilibrium point} if it is stable and attracting. If \(\eta = \infty\), \(x^*\) is said to be \emph{globally asymptotically stable}.
\end{itemize}



\end{document}